% ======================================================================
% SECTION 4: DISCUSSION
% Five thematic subsections expanded from the v0.8.0 bracketed draft at
% 2030-pdac-1min/paper/draft-paper/sections/discussion.tex.
% Location: 2030-pdac-1min/paper/full-paper/sections/discussion.tex
% ======================================================================

\subsection{Significance of the As-Built Artifacts}
\label{sec:discussion-significance}

The five phase Claude Code workflow output paired with this paper
consists of 26 instruction files (v0.5.0) at
\texttt{2030-pdac-1min/paper/instructions/}, the v0.6.0 codegen
tree (15 src modules across Python, C++, and Rust) at
\texttt{2030-pdac-1min/paper/codegen/}, the v0.7.0 execution tree
(16 subdirectories with deterministic run output) at
\texttt{2030-pdac-1min/paper/execution/}, and the v0.8.0 bracketed
draft template at \texttt{2030-pdac-1min/paper/draft-paper/} that
this v0.9.0 full paper expands without modification. The workflow
is end to end in the sense that no human authored a single file
beyond the prompts; the agent stack \cite{claude-code,
claude-opus-47, claude-sonnet-46} carried the five phases.

The significance argument is straightforward: this is the first 60
second 8 arm Whipple paired with a precision oncology drug
development program in the literature.\ The prior 60 second GBM
resection \cite{kawchak_2026_20113157} demonstrated the 4 arm
pattern; the present work scales the pattern to 8 arms and adds
the Daraxonrasib LLM bound advisory layer. The four prior author
PDAC papers \cite{kawchak_2025_15735068, kawchak_2025_16415815,
kawchak_2025_17001137, kawchak_2025_17239510} and the Daraxonrasib
summary \cite{kawchak_2025_18099351} independently complemented
real RASolute 302 trial development throughout 2025; this paper
makes the surgical robot complement explicit and ties the
pharmacologic and surgical workflows under one repository.

\subsection{Opportunistic FDA Framing for Physical AI Oncology Surgery}
\label{sec:discussion-fda}

The 28 April 2026 FDA RTCT announcement \cite{fda2026realtime} is
the regulatory entry point that this work is positioned against.
The two named oncology pharmacology pilots (AstraZeneca TRAVERSE
and Amgen STREAM-SCLC) leave the surgical theater open as a third
pilot family. The present work is positioned as advancing the new
field of Physical AI oncology trial surgical procedures while
keeping the United States Number 1 in the world regarding patient
safety, efficacy, and speed benefits in oncological robotic
clinical trial surgeries. The framing is respectful but
opportunistic.

The four regulatory anchors are the FDA Software as a Medical
Device framework \cite{fda-samd}, the IEC 80601-2-77 robotically
assisted surgical equipment standard \cite{iec-80601-2-77}, the
IEC 62304 medical device software life cycle standard
\cite{iec-62304}, and the 21 CFR 50.30 informed consent obligation
\cite{cfr-21-50-30}. ICH E6(R3) Good Clinical Practice
\cite{ich-e6r3} adds the international clinical trial framing for
the eventual cross border deployment. The reporting floor is set
by TRIPOD+AI \cite{Collins2024TRIPODAI} and CREMLS
\cite{ElEmam2024CREMLS}; both apply to every per round tournament
rationale and to the per iteration advisory output.

The cross jurisdictional opportunity is broader than a single FDA
filing. The EU Medical Device Regulation, the China National
Medical Products Administration, the Japan Pharmaceuticals and
Medical Devices Agency, the United Kingdom Medicines and Healthcare
products Regulatory Agency, Health Canada, and the Therapeutic
Goods Administration of Australia each maintain a parallel SaMD
review track. The on premises LLM thesis preserves the data
residency requirement that governs the parallel non US frameworks
without the need for a separate inference deployment per
jurisdiction. The next subsection details the mechanism
that the SaMD framework will review.

\subsection{On-Premises LLM Single Robot Error Minimization}
\label{sec:discussion-onprem}

The thesis from
\texttt{2030-pdac-1min/paper/instructions/README.md} is stated
verbatim: on premises repository based LLMs provide commands to
standard oncology surgical robots based on real time sensor data
and controlled via x, y, z coordinates to administer patient
treatment; this workflow minimizes single robot error potential.
The mechanism is second opinion oracle isolation. The LLM acts as
an oracle that sits outside the robot vendor kinematic stack, so a
fault in the vendor controller does not propagate into the LLM
advisory layer. The 10 kHz heartbeat bus, the 100 microsecond
watchdog, and the 3 ms E stop budget specified in
\texttt{multi_arm_coordination_8arm.md} and implemented in
\texttt{2030-pdac-1min/paper/codegen/src/coordination/arm_heartbeat_10khz.cpp}
bound every LLM advisory inside the safety floor before it can
reach an actuator.

The quality and safety benefits realized in the v0.7.0 execution
tree are concrete.\ The 5 vessel safety zone gate located at
\texttt{2030-pdac-1min/paper/execution/vascular/gate_verdicts.csv}
logged 100 ticks with the following details: clear 83, no_fly 6, soft_warning 6, and
hard_stop 5; every hard_stop fired the 3 ms E stop within budget.\
The per anastomosis controller outcomes at
\texttt{2030-pdac-1min/paper/execution/anastomosis/} recorded PJ
Grade A in 31 of 32 iterations, HJ leak absent in 14 of 32
iterations per controller (consistent with the bile spectro-
photometry leak detection threshold), and GJ patent in 32 of 32
iterations. The composite mean 93.298 with std 1.225 indicates a
tight outcome envelope around the 32 iteration sample, which is a
quality and safety win against the 2025 Dutch human cohort mean
ideal outcome rate of 47 percent \cite{DutchCohort2025Whipple}
under the cross time normalization caveat.

The privacy and data residency benefits sit on top of the safety
benefits. The on premises deployment never sends patient sensor
data outside the institution. The next subsection
documents how this draft head start accelerates the downstream
final paper pass.

\subsection{Head Start for the Downstream Full Paper Pass}
\label{sec:discussion-headstart}

This v0.9.0 full paper provides the worked example that a
downstream Claude Code Opus 4.7 1M Max pass can reuse on a
sibling cancer site (for example, hepatocellular carcinoma,
cholangiocarcinoma, or gastric cancer) without re authoring the
section structure, the table column type contract, the reference
inventory pattern, or the cross study cite framing. The bracketed
draft template at
\texttt{2030-pdac-1min/paper/draft-paper/} (preserved verbatim in
this PR) names the exact source files in
\texttt{2030-pdac-1min/paper/\{instructions,codegen,execution,inputs\}/} that the downstream pass should process; this v0.9.0
full paper resolves those brackets into running prose, anchored
tables, and ASCII diagrams without inventing new sources.

The benefits of the head start are tangible. The downstream pass
can focus its 1M context budget on synthesis rather than
discovery; the file name level specificity in each bracket
reduces hallucination risk; the pre populated tables anchor the
dimensional and unit invariants; the references.bib starter ships
with the doi plus url plus note triad that the downstream pass
extends rather than replaces. The v0.4.0 GBM paper at
\cite{kawchak_2026_20113157} provides the parallel cross study
anchor; the v0.9.0 PDAC paper provides the second worked example
that closes the GBM to PDAC pattern.\ A future sibling
hepatocellular carcinoma paper would start by copying the v0.8.0
PDAC draft template, swapping the disease specific instruction
files, and resolving the brackets in one pass.

\subsection{Practical Real-Life Adoption Insights}
\label{sec:discussion-realworld}

The honest accounting of how close this work sits to real life is
the most important paragraph in the discussion. The 60 second
simulation is simulation against simulation for the three
hypothetical 2030 robot competitors and simulation against the
2025 Dutch cohort numbers \cite{DutchCohort2025Whipple} for the
human baseline.\ The synthetic patient PAT-PDAC-0001 from
\texttt{2030-pdac-1min/paper/instructions/pdac_context_1min.md}
is one synthetic patient, not a cohort. The TRIPOD+AI
\cite{Collins2024TRIPODAI} and CREMLS \cite{ElEmam2024CREMLS}
reporting standards must be satisfied before any clinical
translation step. The 2025 SOTA platforms cited in
\cite{intuitive-davinci-sp, medtronic-hugo-ras}
cannot perform a 60 second Whipple today; bench top prototypes
that close the 5x to 500x performance gap are a multi year
hardware program with FDA Breakthrough Device pathway gating.

The Daraxonrasib bridge is closer to real life because the drug
is real. Daraxonrasib holds an FDA Breakthrough Therapy
designation from June 2025 \cite{rev-fda-breakthrough}, is
enrolled in the RASolute 302 phase 3 trial in previously treated
metastatic pancreatic cancer \cite{rasolute302}, and is enrolled
in the RASolve 301 first line trial with pembrolizumab
\cite{rasolve301}.\ The perioperative pause and restart logic in
\texttt{2030-pdac-1min/paper/instructions/daraxonrasib_integration.md}
is consistent with the public protocols. The LLM bound advisory
layer is the new piece this work contributes; an FDA Q-Sub pre
submission under the SaMD framework \cite{fda-samd} is the
tractable next regulatory step. Table \ref{tab:discussion-realworld}
summarizes the 5 real life adoption gaps and the tractable next
step for each.

\begin{table}[H]
\caption{5 real life adoption gaps and tractable next steps for the 60 second 8 arm Whipple plus Daraxonrasib pairing.}
\label{tab:discussion-realworld}
\centering
\begin{tabular}{>{\raggedright\arraybackslash}p{4.8cm} >{\raggedright\arraybackslash}p{5.4cm} >{\raggedright\arraybackslash}p{5.2cm}}
\toprule
\textbf{Gap} & \textbf{Why it matters} & \textbf{Tractable next step} \\
\midrule
  Synthetic patient PAT-PDAC-0001 & Single patient, no cohort variance & Cohort scale up under TRIPOD+AI \cite{Collins2024TRIPODAI} \\
  Hypothetical PancreSpeed 1.0 hardware & 5x to 500x gap vs 2025 SOTA & Multi year hardware program with FDA Breakthrough \\
  Simulation vs simulation for 3 robot competitors & No real bench data on competitors & Vendor handoff under cross licensing terms \\
  60 second time scale not biologically validated & Tissue creep, anastomosis healing, pharmacokinetics & Translational work on porcine, human cadaver models \\
  Not yet IRB approved & SaMD review pending & Pre submission Q-Sub \cite{fda-samd} \\
\bottomrule
\end{tabular}
\end{table}
